\documentclass[11pt]{article}
\usepackage[a4paper,margin=1in]{geometry}
\usepackage{amsmath,amssymb,mathtools,booktabs,enumitem}
\usepackage[dvipsnames]{xcolor}
\usepackage{hyperref,microtype,comment,pgfplots}
\pgfplotsset{compat=1.18}
\hypersetup{colorlinks=true,linkcolor=MidnightBlue,urlcolor=MidnightBlue,
 pdftitle={Recitation 10: Power and the Sign Test},pdfauthor={Xuanxi Zhang}}
% Default: student handout. Define \WITHSOLUTIONS for answers and pacing.
% Both versions compile from this self-contained file; no external figures.
\newif\ifsolutions
\ifdefined\WITHSOLUTIONS\solutionstrue\else\solutionsfalse\fi
\newcommand{\rv}[1]{\widetilde{#1}}
\newcommand{\E}{\mathbb E}
\newcommand{\Pbb}{\mathbb P}
\newcommand{\Var}{\operatorname{Var}}
\newcommand{\pow}{\operatorname{pow}}
\ifsolutions
 \newenvironment{answer}{\par\medskip\noindent\textcolor{MidnightBlue}{\textbf{Solution.}}\ }{\par\medskip}
 \newcommand{\workingspace}[1]{}
\else
 \excludecomment{answer}
 \newcommand{\workingspace}[1]{\vspace{#1}}
\fi
\setlist[enumerate]{label=(\alph*),leftmargin=2em,itemsep=0.4em}
\setlist[itemize]{leftmargin=1.5em,itemsep=0.3em}
\title{\huge Recitation 10\\[0.35em]\normalsize Power and the Sign Test\\[0.35em]
 \large\ifsolutions Solutions and Teaching Notes\else Review and Practice Problems\fi}
\author{Xuanxi Zhang}
\date{November 13, 2026}
\begin{document}
\maketitle
\noindent\textit{Reading:}
\href{https://github.com/cfgranda/ps4ds/blob/main/ps4ds_preprint.pdf}
{\emph{Probability and Statistics for Data Science}},
Section 10.6, Example 10.17, and Exercise 10.10.

\section*{Review: how often will a test detect an effect?}
For a fixed test and significance level $\alpha$, its \textbf{power function} is
\[
 \pow(\theta)=\Pbb_\theta(\text{reject }H_0).
\]
Under an alternative, this is the probability of a true positive;
$1-\pow(\theta)$ is the probability of a type II error.
Under a null parameter, the same rejection probability is the type I error
rate. A level-$\alpha$ test bounds it by $\alpha$ throughout the null.
A $p$ value evaluates the observed data under the null; power evaluates a
\emph{decision rule} across repeated datasets at a specified parameter.

\paragraph{1. Specify the free-throw model (Example 10.17).}
In the free-throw example, Antetokounmpo made 34 of 44 home attempts and
22 of 41 away attempts. For repeated datasets of these sizes, assume all shots are
independent, with fixed success probabilities $\theta_H$ and $\theta_A$:
\[
 \rv h_i\sim\operatorname{Bernoulli}(\theta_H),\qquad
 \rv a_j\sim\operatorname{Bernoulli}(\theta_A),\qquad
 \rv d=\frac1{44}\sum_{i=1}^{44}\rv h_i-
       \frac1{41}\sum_{j=1}^{41}\rv a_j.
\]
The observed difference is $d_{\mathrm{obs}}=34/44-22/41\approx0.236$.
The right-tailed test uses $H_0:\theta_H\leq\theta_A$ versus
$H_1:\theta_H>\theta_A$; the two-tailed test uses
$H_0:\theta_H=\theta_A$ versus $H_1:\theta_H\ne\theta_A$.
Both are calibrated at equality.

\paragraph{2. Use the null distribution to set the rejection region.}
Follow the book's power approximation: fix $\theta_H=0.685$ and use the
\emph{true null} standard deviation at $\theta_A=\theta_H$,
\[
 \sigma_{\mathrm{null}}=
 \sqrt{0.685(1-0.685)\left(\frac1{44}+\frac1{41}\right)}
 \approx0.10083.
\]
Write $\Phi(u)=\Pbb(\rv z\leq u)$ for $\rv z\sim\mathcal N(0,1)$.
At $\alpha=0.05$, solving the normal-tail equations gives
\[
 \begin{array}{lll}
 \text{One tail:}&1-\Phi(\tau_1/\sigma_{\mathrm{null}})=\alpha,
   &\tau_1=\sigma_{\mathrm{null}}\Phi^{-1}(0.95)\approx0.16585,\\[0.5em]
 \text{Two tails:}&2[1-\Phi(\tau_2/\sigma_{\mathrm{null}})]=\alpha,
   &\tau_2=\sigma_{\mathrm{null}}\Phi^{-1}(0.975)\approx0.19762.
 \end{array}
\]
Reject when $d_{\mathrm{obs}}\geq\tau_1$ or $|d_{\mathrm{obs}}|\geq\tau_2$,
respectively. \textbf{Keep these cutoffs fixed as $\theta_A$ varies.}
The observed pooled SE, about $0.10291$, is a different quantity; the
calculation here follows the fixed-null approximation in equation (10.82).

\newpage
\section*{Review: evaluate the rejection region under an alternative}
\paragraph{3. Change the distribution, keeping the test fixed.}
At a specified pair $(\theta_H,\theta_A)$, let $\delta=\theta_H-\theta_A$.
Independence and the Gaussian approximation to the sample proportions give
\[
 \rv d\ \approx\ \mathcal N(\delta,\sigma_{\mathrm{pow}}^2),\qquad
 \sigma_{\mathrm{pow}}^2=
 \frac{\theta_H(1-\theta_H)}{44}+\frac{\theta_A(1-\theta_A)}{41}.
\]
\textbf{The null SD sets the cutoffs; the alternative SD computes power.}
Standardize under this specified model, subtracting its mean $\delta$:
\[
 \boxed{\pow_1(\theta_H,\theta_A)
 \approx1-\Phi\!\left(\frac{\tau_1-\delta}{\sigma_{\mathrm{pow}}}\right).}
\]
For a two-tailed test, add the probabilities of the two rejection tails:
\[
 \boxed{\pow_2(\theta_H,\theta_A)\approx
 \Phi\!\left(\frac{-\tau_2-\delta}{\sigma_{\mathrm{pow}}}\right)
 +1-\Phi\!\left(\frac{\tau_2-\delta}{\sigma_{\mathrm{pow}}}\right).}
\]

\paragraph{4. Work one example: $\theta_H=0.685$, $\theta_A=0.55$.}
Here $\delta=0.135$ and $\sigma_{\mathrm{pow}}\approx0.10460$. Thus
\[
 \pow_1\approx1-\Phi(0.29495)=0.3840,\qquad
 \pow_2\approx\Phi(-3.18005)+1-\Phi(0.59872)=0.2754.
\]
The right-tailed test detects this specified difference about 38.4\% of the
time in this approximation; it misses it about 61.6\% of the time.
This is not a probability that the alternative is true.

\begin{center}
\begin{tikzpicture}
\begin{axis}[
 width=0.91\textwidth,height=5.25cm,
 xmin=0,xmax=1,ymin=0,ymax=1.03,
 xlabel={Away success probability $\theta_A$},ylabel={Rejection probability},
 xtick={0,0.25,0.5,0.685,0.85,1},
 xticklabels={$0$,$0.25$,$0.5$,$0.685$,$0.85$,$1$},
 ytick={0,0.25,0.5,0.75,1},tick label style={font=\small},
 label style={font=\small},
 legend style={at={(0.03,0.05)},anchor=south west,font=\small,
   draw=none,fill=white},
 grid=major,grid style={gray!18}]
% Numeric coordinates are computed from the displayed Gaussian formulas.
\addplot[black,thick] coordinates {
(0.000,1.0000000) (0.010,1.0000000) (0.020,1.0000000) (0.030,1.0000000) (0.040,1.0000000) (0.050,1.0000000) (0.060,1.0000000) (0.070,1.0000000) (0.080,1.0000000) (0.090,0.9999999) (0.100,0.9999997) (0.110,0.9999992) (0.120,0.9999980) (0.130,0.9999956) (0.140,0.9999907) (0.150,0.9999814) (0.160,0.9999641) (0.170,0.9999338) (0.180,0.9998823) (0.190,0.9997980) (0.200,0.9996642) (0.210,0.9994580) (0.220,0.9991488) (0.230,0.9986969) (0.240,0.9980520) (0.250,0.9971515) (0.260,0.9959201) (0.270,0.9942684) (0.280,0.9920929) (0.290,0.9892765) (0.300,0.9856889) (0.310,0.9811886) (0.320,0.9756252) (0.330,0.9688419) (0.340,0.9606798) (0.350,0.9509817) (0.360,0.9395968) (0.370,0.9263854) (0.380,0.9112246) (0.390,0.8940125) (0.400,0.8746733) (0.410,0.8531616) (0.420,0.8294661) (0.430,0.8036120) (0.440,0.7756636) (0.450,0.7457243) (0.460,0.7139371) (0.470,0.6804828) (0.480,0.6455773) (0.490,0.6094680) (0.500,0.5724293) (0.510,0.5347567) (0.520,0.4967605) (0.530,0.4587594) (0.540,0.4210729) (0.550,0.3840145) (0.560,0.3478843) (0.570,0.3129628) (0.580,0.2795044) (0.590,0.2477327) (0.600,0.2178359) (0.610,0.1899637) (0.620,0.1642256) (0.630,0.1406898) (0.640,0.1193841) (0.650,0.1002970) (0.660,0.0833813) (0.670,0.0685571) (0.680,0.0557165) (0.685,0.0500000) (0.690,0.0447288) (0.700,0.0354456) (0.710,0.0277067) (0.720,0.0213450) (0.730,0.0161925) (0.740,0.0120840) (0.750,0.0088619) (0.760,0.0063789) (0.770,0.0045011) (0.780,0.0031089) (0.790,0.0020987) (0.800,0.0013822) (0.810,0.0008864) (0.820,0.0005523) (0.830,0.0003336) (0.840,0.0001947) (0.850,0.0001096) (0.860,0.0000592) (0.870,0.0000306) (0.880,0.0000151) (0.890,0.0000070) (0.900,0.0000031) (0.910,0.0000013) (0.920,0.0000005) (0.930,0.0000002) (0.940,0.0000001) (0.950,0.0000000) (0.960,0.0000000) (0.970,0.0000000) (0.980,0.0000000) (0.990,0.0000000) (1.000,0.0000000)
};
\addlegendentry{One tail}
\addplot[MidnightBlue,thick,densely dotted] coordinates {
(0.000,1.0000000) (0.010,1.0000000) (0.020,1.0000000) (0.030,1.0000000) (0.040,1.0000000) (0.050,1.0000000) (0.060,1.0000000) (0.070,0.9999999) (0.080,0.9999997) (0.090,0.9999991) (0.100,0.9999979) (0.110,0.9999950) (0.120,0.9999892) (0.130,0.9999777) (0.140,0.9999563) (0.150,0.9999180) (0.160,0.9998523) (0.170,0.9997438) (0.180,0.9995707) (0.190,0.9993033) (0.200,0.9989019) (0.210,0.9983155) (0.220,0.9974803) (0.230,0.9963179) (0.240,0.9947350) (0.250,0.9926230) (0.260,0.9898579) (0.270,0.9863022) (0.280,0.9818061) (0.290,0.9762105) (0.300,0.9693507) (0.310,0.9610599) (0.320,0.9511747) (0.330,0.9395399) (0.340,0.9260138) (0.350,0.9104738) (0.360,0.8928217) (0.370,0.8729882) (0.380,0.8509376) (0.390,0.8266707) (0.400,0.8002277) (0.410,0.7716894) (0.420,0.7411773) (0.430,0.7088527) (0.440,0.6749146) (0.450,0.6395961) (0.460,0.6031602) (0.470,0.5658945) (0.480,0.5281051) (0.490,0.4901099) (0.500,0.4522322) (0.510,0.4147931) (0.520,0.3781053) (0.530,0.3424663) (0.540,0.3081532) (0.550,0.2754171) (0.560,0.2444797) (0.570,0.2155301) (0.580,0.1887234) (0.590,0.1641801) (0.600,0.1419871) (0.610,0.1221994) (0.620,0.1048436) (0.630,0.0899213) (0.640,0.0774145) (0.650,0.0672902) (0.660,0.0595068) (0.670,0.0540191) (0.680,0.0507842) (0.685,0.0500000) (0.690,0.0497664) (0.700,0.0509417) (0.710,0.0543012) (0.720,0.0598538) (0.730,0.0676270) (0.740,0.0776670) (0.750,0.0900372) (0.760,0.1048146) (0.770,0.1220861) (0.780,0.1419417) (0.790,0.1644678) (0.800,0.1897386) (0.810,0.2178069) (0.820,0.2486939) (0.830,0.2823791) (0.840,0.3187899) (0.850,0.3577919) (0.860,0.3991799) (0.870,0.4426711) (0.880,0.4879002) (0.890,0.5344180) (0.900,0.5816938) (0.910,0.6291229) (0.920,0.6760389) (0.930,0.7217332) (0.940,0.7654802) (0.950,0.8065693) (0.960,0.8443413) (0.970,0.8782293) (0.980,0.9078001) (0.990,0.9327916) (1.000,0.9531425)
};
\addlegendentry{Two tails}
\addplot[gray,dashed,domain=0:1] {0.05};
\addlegendentry{$\alpha=0.05$}
\addplot[gray,densely dashed,forget plot] coordinates {(0.685,0) (0.685,1.03)};
\end{axis}
\end{tikzpicture}
\end{center}
\noindent\small
Both curves equal $0.05$ at equality in this Gaussian approximation.
One tail gives greater power when home is better; two tails also detect
when away is better. Choose the direction before observing outcomes.
The actual pooled test re-estimates its null SE for every dataset; its
power can be computed by simulation or by summing binomial probabilities.
\normalsize

\newpage
\section*{Practice problems}
\noindent Three core problems; Problem 4 is a backup. The indicated times
are for independent work before discussion.

\subsection*{Problem 1: Which effects can the free-throw test detect?}
\noindent\textit{Independent work: 8 minutes. Example 10.17 and Exercise 10.5.}
Let $\theta_H,\theta_A$ denote the home/away success probabilities.
Assume all 44 home shots and 41 away shots are independent, and fix
$\theta_H=0.685$. Let $d$ be the observed home-minus-away proportion
and $\rv d$ its random counterpart. Compare these tests:
\begin{center}
\begin{tabular}{@{}llll@{}}
\toprule Test & $H_0$ & $H_1$ & Reject when\\\midrule
One tail & $\theta_H\leq\theta_A$ & $\theta_H>\theta_A$ & $d\geq0.16585$\\
Two tails & $\theta_H=\theta_A$ & $\theta_H\ne\theta_A$ & $|d|\geq0.19762$\\
\bottomrule
\end{tabular}
\end{center}
These are the fixed-null Gaussian cutoffs for nominal $\alpha=0.05$,
calibrated at $\theta_H=\theta_A=0.685$ as on page 1.
\begin{enumerate}
 \item Derive $\E[\rv d]=\theta_H-\theta_A$ and
 $\Var(\rv d)=\theta_H(1-\theta_H)/44+\theta_A(1-\theta_A)/41$.
 Why do the two variance contributions add?
 \item Suppose $\theta_H=0.685$ and $\theta_A=0.75$.
 Compute $\delta=\theta_H-\theta_A$, $\sigma_{\mathrm{pow}}$, and
 the approximate rejection probability of each test.
 You may use $\Phi(2.371)\approx0.9911$,
 $\Phi(1.362)\approx0.9134$, and $\Phi(2.698)\approx0.9965$.
 \item For each test, is the scenario in (b) in $H_0$ or $H_1$?
 Explain the small one-tailed rejection probability. Which test detects
 differences in either direction? At $\theta_A=0.685$, what is each
 approximate rejection probability, and which type of error does it represent?
\end{enumerate}
\workingspace{9cm}
\begin{answer}
\textbf{(a)} A mean of $n$ independent Bernoulli$(\theta)$ observations
has expectation $\theta$ and variance $\theta(1-\theta)/n$.
Linearity of expectation and independence between groups therefore give
\[
 \E[\rv d]=\theta_H-\theta_A,\qquad
 \Var(\rv d)=\frac{\theta_H(1-\theta_H)}{44}
                         +\frac{\theta_A(1-\theta_A)}{41}.
\]
The variances add because the covariance is zero and the coefficient
$-1$ is squared when computing variance.

\textbf{(b)} Now $\delta=-0.065$ and
\[
 \sigma_{\mathrm{pow}}=
 \sqrt{\frac{0.685(0.315)}{44}+\frac{0.75(0.25)}{41}}
 \approx0.09735.
\]
Using $\tau_1\approx0.16585$ and $\tau_2\approx0.19762$,
\begin{align*}
 \pow_1&\approx1-\Phi\!\left(\frac{0.16585+0.065}{0.09735}\right)
       \approx\boxed{0.0089},\\
 \pow_2&\approx\Phi\!\left(\frac{-0.19762+0.065}{0.09735}\right)
       +1-\Phi\!\left(\frac{0.19762+0.065}{0.09735}\right)
       \approx\boxed{0.0900}.
\end{align*}
The supplied rounded cdf values give about $0.0901$ for the second result.

\textbf{(c)} The right-tailed test rejects for a sufficiently positive
home-minus-away difference, but the true difference here is negative.
The two-tailed test can detect either direction, although a 6.5-point gap
is still hard to detect with these sample sizes. In part (b),
$0.685<0.75$: the one-tailed $H_0$ holds, but the two-tailed $H_1$ holds.
Thus $0.0089$ is a null rejection probability, whereas $0.0900$ is power
against the two-tailed alternative.
At equality, $\delta=0$ and $\sigma_{\mathrm{pow}}=\sigma_{\mathrm{null}}$,
so the two Gaussian formulas both give $\boxed{0.05}$. Both nulls hold
there, so these are type I error probabilities in the approximation.
\end{answer}

\newpage
\subsection*{Problem 2: Derive the sign test (Exercise 10.10(a)--(b))}
\noindent\textit{Independent work: 7 minutes.}
For paired measurements $(x_i,y_i)$, the sign statistic counts how often
the second measurement is larger:
\[
 t=\sum_{i=1}^n\mathbf 1(y_i>x_i).
\]
Its random counterpart is $\rv t=\sum_{i=1}^n\mathbf 1(\rv y_i>\rv x_i)$.
Assume independent pairs, no ties, and a common sign probability
$q=\Pbb(\rv y_i>\rv x_i)$. The one-sided sign test uses
\[
 \boxed{H_0:q\leq\tfrac12\qquad\text{versus}\qquad H_1:q>\tfrac12.}
\]
Large values of $t$ support $H_1$. Calibrate the test at the null boundary
$q=1/2$; parts (a)--(b) derive this case from a same-distribution model.
\begin{enumerate}
 \item For independent $\rv a,\rv b$ with a common continuous distribution,
 show that \mbox{$\Pbb(\rv a>\rv b)=1/2$}.
 Where do continuity and symmetry enter?
 \item Suppose all $2n$ measurements are independent samples from the
 same continuous distribution. Show that this gives the null boundary
 $q=1/2$, and derive the distribution of $\rv t$ and its pmf.
 State the upper-tail $p$-value formula for the stated alternative.
\end{enumerate}
\noindent\textit{Discussion:} Why does this test not require a Gaussian
model for the measurements? What information does it discard?
\workingspace{10cm}
\begin{answer}
\textbf{(a)} Independence and the common distribution make
$(\rv a,\rv b)$ and $(\rv b,\rv a)$ identically distributed. Therefore
$\Pbb(\rv a>\rv b)=\Pbb(\rv b>\rv a)$.
Independence and continuity imply $\Pbb(\rv a=\rv b)=0$. The two strict
inequality events partition the sample space up to this zero-probability
event, so each has probability $\boxed{1/2}$.

\textbf{(b)} Define $\rv s_i=\mathbf 1(\rv y_i>\rv x_i)$.
Part (a) gives $\rv s_i\sim\operatorname{Bernoulli}(1/2)$.
Because distinct pairs are independent, their signs are independent.
Consequently,
\[
 \boxed{\rv t\sim\operatorname{Binomial}(n,1/2),\qquad
 p_{\rv t}(k)=\binom nk2^{-n},\quad k=0,\ldots,n.}
\]
Large counts support the stated alternative. For observed count $t$,
\[
 \boxed{p_{\mathrm{right}}(t)=\Pbb_0(\rv t\geq t)
       =\sum_{k=t}^{n}\binom nk2^{-n}.}
\]
Include the observed count in this tail: this is a discrete test.
For $q\leq1/2$, the binomial upper tail is largest at $q=1/2$, so this
boundary calibration controls error throughout the stated null.
The continuous measurement distribution need not be specified: at the
null boundary, the independent signs are Bernoulli$(1/2)$.
The statistic uses only which measurement is larger, discarding the
magnitudes of the differences.

\medskip
\noindent\textit{Assumption check:} Independent pairs alone do not imply
independent measurements within a pair. For paired data, within-pair
\emph{exchangeability} (swapping the labels leaves the joint distribution
unchanged) and no ties also give fair signs, even if the two measurements
are dependent. Equal marginal distributions alone are insufficient.
\end{answer}

\newpage
\subsection*{Problem 3: Are left ears longer? (Exercise 10.10(c)--(d))}
\noindent\textit{Independent work: 10 minutes.}
A friend claims that left ears tend to be longer than right ears.
Measurements from ten people, in inches, are
\begin{center}
\begin{tabular}{@{}l*{10}{r}@{}}
\toprule
Left  &2.4&2.7&3.2&2.3&2.0&2.6&3.2&2.3&2.9&2.3\\
Right &2.2&2.6&3.3&2.2&2.1&2.5&3.1&2.5&2.7&2.2\\
\bottomrule
\end{tabular}
\end{center}
Assume independent people and no ties. Set $y_i$ to the left-ear length,
$x_i$ to the right-ear length, and let $q=\Pbb(\rv y_i>\rv x_i)$ be common
across people. Test at significance level $\alpha=0.05$:
\[
 \boxed{H_0:q\leq\tfrac12\qquad\text{versus}\qquad H_1:q>\tfrac12.}
\]
Use $t=\sum_{i=1}^{10}\mathbf 1(y_i>x_i)$; large counts support $H_1$.
Under the null boundary $q=1/2$, $\rv t\sim\operatorname{Binomial}(10,1/2)$.
\begin{enumerate}
 \item Explain $H_1$ in words, list the ten signs of $y_i-x_i$, and compute $t$.
 \item Compute the exact one-sided $p$ value and give the decision at
 $\alpha=0.05$. Does the decision establish equal ear lengths?
 \item As an extension, find the smallest count that rejects at 5\%.
 What is the rejection probability at $q=1/2$? Is it exactly 5\%?
 \item Return to Exercise 10.10(d): suppose all measurements are
 independent, and $\rv x_i$ and $\rv y_i$ each have the same continuous
 distribution $F_i$ within each pair, but $F_i$ may differ across people.
 Is the same sign test still valid? Justify your answer.
\end{enumerate}
\noindent\textit{Arithmetic aid:}
$\binom{10}{7}=120$, $\binom{10}{8}=45$,
$\binom{10}{9}=10$, $\binom{10}{10}=1$, and $2^{10}=1024$.
\workingspace{8cm}
\begin{answer}
\textbf{(a)} The alternative $H_1:q>1/2$ means that a person's left ear
is more likely to be longer than the right ear. The signs of
left minus right are
\[
 (+,+,-,+,-,+,+,-,+,+),\qquad \boxed{t=7}.
\]
The null $q\leq1/2$ says that left-longer is no more likely than right-longer;
it does not assert exact equality of every person's two ear lengths.

\textbf{(b)} At the null boundary $q=1/2$, $\rv t\sim\operatorname{Binomial}(10,1/2)$.
The exact tail probability is
\[
 p=\Pbb_0(\rv t\geq7)
 =\frac{120+45+10+1}{1024}
 =\boxed{0.171875}.
\]
Since $p>0.05$, do not reject. These data do not provide sufficient
evidence for the directional claim at 5\%; they do not prove equal
lengths or equal distributions. A prespecified two-sided test would give
$2(0.171875)=0.34375$.

\textbf{(c)} The two relevant tails are
\[
 \Pbb_0(\rv t\geq8)=\frac{56}{1024}=0.0546875>0.05,
 \qquad
 \Pbb_0(\rv t\geq9)=\frac{11}{1024}\approx0.0107422<0.05.
\]
Thus reject exactly when $\boxed{t\geq9}$.
At $q=1/2$, the rejection probability is about $1.07\%$, and it is smaller
for $q<1/2$. Discreteness makes this nonrandomized level-$0.05$ test conservative.

\textbf{(d)} Yes. For each $i$, independence, continuity, and the common
within-pair distribution $F_i$ give
$\Pbb(\rv y_i>\rv x_i)=1/2$. Independent pairs therefore still yield
i.i.d. Bernoulli$(1/2)$ signs, even though the underlying lengths are
not identically distributed across people. Their sum remains
$\boxed{\operatorname{Binomial}(10,1/2)}$.
\end{answer}

\newpage
\subsection*{Problem 4 (backup): Gaussian power and sample size}
\noindent\textit{Independent work: 5 minutes. Archived Recitation 7,
with a power extension.}
Let $\rv x_1,\ldots,\rv x_n$ be i.i.d. $\mathcal N(\theta,1)$.
We test $H_0:\theta\leq0$ against $H_1:\theta>0$ using the sample mean.
Let $m(X)=n^{-1}\sum_i x_i$ denote its observed value.
\begin{enumerate}
 \item Construct a level-$0.05$ test. Explain why checking $\theta=0$
 suffices to bound the rejection probability for every $\theta\leq0$.
 \item Derive the power function at any fixed $\theta$.
 \item How large must $n$ be for power at least $0.80$ when
 $\theta=0.5$? Use $\Phi^{-1}(0.95)=1.64485$ and
 $\Phi^{-1}(0.80)=0.84162$.
\end{enumerate}
\workingspace{11cm}
\begin{answer}
\textbf{(a)} The random sample mean has distribution
$\rv m\sim\mathcal N(\theta,1/n)$, so at $\theta=0$,
$\sqrt n\,\rv m\sim\mathcal N(0,1)$. Reject if
\[
 \boxed{m(X)\geq\frac{\Phi^{-1}(0.95)}{\sqrt n}.}
\]
For $\theta<0$, the distribution shifts left while the cutoff stays
fixed, so its upper-tail rejection probability decreases. The maximum
null rejection probability is $0.05$, attained at $\theta=0$.

\textbf{(b)} Standardizing under the true parameter gives the exact power
function (the observations are Gaussian):
\[
 \pow(\theta)=\Pbb_\theta\!\left(\rv m\geq
                          \frac{\Phi^{-1}(0.95)}{\sqrt n}\right)
 =\boxed{1-\Phi\!\left(\Phi^{-1}(0.95)-\sqrt n\,\theta\right).}
\]
For fixed $\theta>0$, increasing $n$ increases power.

\textbf{(c)} Require
$\Phi(1.64485-0.5\sqrt n)\leq0.20$. Since
$\Phi^{-1}(0.20)=-0.84162$,
\[
 \sqrt n\geq\frac{1.64485+0.84162}{0.5},\qquad
 n\geq24.7301.
\]
Round up to get $\boxed{n=25}$.
The resulting power is about $0.8038$.
\end{answer}

\ifsolutions
\newpage
\section*{Instructor notes: 75-minute plan}
\begin{center}
\begin{tabular}{@{}p{0.16\textwidth}p{0.77\textwidth}@{}}
\toprule Time & Activity\\\midrule
0--5 min & Define power, type I and type II errors; distinguish power
from a $p$ value.\\
5--13 min & Recall the free throws. Fix the null SD and derive the
one- and two-tailed cutoffs.\\
13--22 min & Shift to the specified alternative, derive both powers,
and work $\theta_A=0.55$. Read the power curves.\\
22--25 min & Introduce paired signs: a fair sign is Bernoulli$(1/2)$.
Explain the setup for Exercise 10.10 without giving the proof.\\
25--41 min & Problem 1: 8 minutes independent work, 8 minutes discussion.\\
41--55 min & Problem 2: 7 minutes independent work, 7 minutes discussion.\\
55--72 min & Problem 3: 10 minutes independent work, 7 minutes discussion.\\
72--75 min & Exit check and questions.\\\bottomrule
\end{tabular}
\end{center}
\noindent\textit{Fast-class option:} Keep independent-work times intact.
Shorten the three core discussions to 6, 4, and 4 minutes, respectively.
This finishes Problem 3 at minute 64. Use 64--69 for independent work on
Problem 4, 69--72 for its discussion, and 72--75 for the exit check.

\paragraph{Board sequence.}
Write
\[
 \text{null model}\ \longrightarrow\ \text{fixed rejection region}
 \ \longrightarrow\ \text{probability under a specified model}.
\]
Sketch a null bell curve centered at zero and an alternative curve centered
at $0.135$, with the same vertical cutoff $0.16585$ on both.
The area to its right is about $0.05$ under the null and $0.384$ under
the specified alternative. For two tails, draw both $\pm0.19762$ cutoffs
and shade the two alternative tail areas. Then move to the binomial
sign statistic: the same rejection-probability idea applies to discrete data.

\paragraph{Exit check.}
Ask: ``Which SD sets the free-throw cutoff, and which SD computes power?
Why can a level-$0.05$ sign test reject less than 5\% of the time under
the null?'' Expected answers: the null SD and specified-model SD,
respectively; discrete tail probabilities need not equal the target level.

\paragraph{Paired-data assumptions.}
Real ears from the same person can be dependent. The exercise's independent
continuous measurements are sufficient for its derivation; independent
people with exchangeable left/right measurements and no ties are also
sufficient. Do not infer exchangeability from equal marginals alone.
If ties occur, the simple no-ties derivation does not apply directly.
The usual sign test drops ties and uses the number of non-tied pairs;
this requires fair signs conditional on being non-tied. There are no ties
in the supplied data.

\newpage
\section*{Instructor notes: the actual pooled test}
\paragraph{Why the two power calculations differ.}
Example 10.17 fixes the null SD at the true equality model with
$\theta_H=0.685$. In practice, the two-sample $z$ test estimates a pooled
null probability from \emph{each new dataset}. To evaluate that actual
procedure at specified $(\theta_H,\theta_A)$, repeat the following:
\begin{enumerate}[label=\arabic*.]
 \item Simulate independent counts
 $\rv k_H\sim\operatorname{Binomial}(44,\theta_H)$ and
 $\rv k_A\sim\operatorname{Binomial}(41,\theta_A)$.
 \item For the simulated counts $k_H,k_A$, compute
 \[
 d=\frac{k_H}{44}-\frac{k_A}{41},\qquad
 \theta_{\mathrm{ML},0}=\frac{k_H+k_A}{85},\qquad
 s_0=\sqrt{\theta_{\mathrm{ML},0}(1-\theta_{\mathrm{ML},0})
                  (1/44+1/41)}.
 \]
 \item Compute $z=d/s_0$ and use
 $p_1=1-\Phi(z)$ or $p_2=2[1-\Phi(|z|)]$.
 If all attempts are makes or all are misses, $d=s_0=0$;
 define the procedure to not reject (set $p=1$).
 \item Record $\mathbf 1(p\leq0.05)$.
 The fraction of rejections across $B$ independent repetitions estimates
 power. Its estimated Monte Carlo SE is
 $\sqrt{p_{\mathrm{MC}}(1-p_{\mathrm{MC}})/B}$,
 where $p_{\mathrm{MC}}$ denotes that observed fraction.
\end{enumerate}
Because there are only $45\times42$ possible count pairs, we can also
compute the rejection probability directly, without simulation error:
\[
 \sum_{k_H=0}^{44}\sum_{k_A=0}^{41}
 \mathbf 1\{\text{pooled rule rejects at }(k_H,k_A)\}
 \binom{44}{k_H}\theta_H^{k_H}(1-\theta_H)^{44-k_H}
 \binom{41}{k_A}\theta_A^{k_A}(1-\theta_A)^{41-k_A}.
\]
\begin{center}
\begin{tabular}{@{}lrrrr@{}}
\toprule
 &\multicolumn{2}{c}{Fixed-null Gaussian approximation}
 &\multicolumn{2}{c}{Pooled rule: exact summation}\\
\cmidrule(lr){2-3}\cmidrule(l){4-5}
$\theta_A$ & One tail & Two tails & One tail & Two tails\\\midrule
0.550 &0.3840&0.2754&0.3390&0.2574\\
0.685 &0.0500&0.0500&0.0491&0.0496\\
0.750 &0.0089&0.0900&0.0112&0.0998\\
\bottomrule
\end{tabular}
\end{center}
\noindent Here $\theta_H=0.685$ throughout. Exact summation evaluates the
chosen pooled $z$ rule under the independent-binomial model; it does not
turn that rule into an exact level-$0.05$ test for all null parameters.

\paragraph{Hypothesis conventions.}
The book computes the free-throw curves using the equality null.
This handout makes the directional formulation explicit:
$H_0:\theta_H\leq\theta_A$ versus $H_1:\theta_H>\theta_A$, calibrated
at the same equality boundary. Its rejection curve is unchanged; the
away-better region is now part of the directional null. For the sign test,
independent continuous measurements with the same distribution within each
pair imply $q=1/2$, the boundary of our directional null $q\leq1/2$.

\paragraph{Sources and numerical clarification.}
The review and Problem 1 follow
\href{https://github.com/cfgranda/ps4ds/blob/main/ps4ds_preprint.pdf}
{Example 10.17 and Exercise 10.5}; Problems 2--3 cover all of Exercise
10.10, with an added discrete-cutoff question. In the preprint, the
quoted cutoffs $0.166$ and $0.198$ agree with the fixed true-null SD
$0.10083$, not the stated observed SD $0.103$. Using the actual observed
pooled SE $0.10291$ instead gives $0.16927$ and $0.20170$.
The power curves here are recomputed from equations (10.79)--(10.82).
Problem 4(a) reuses the Gaussian test in the Fall 2025 archive's
\texttt{mine/r7/r7.tex}, Section 2; parts (b)--(c) add power and sample-size
planning. The archive had no close match for the requested sign exercise.
\fi
\end{document}
